\documentclass[12pt,a4paper]{article}
% ============================================================
% Preámbulo común a todos los documentos del Proyecto 2
% ============================================================
\usepackage[utf8]{inputenc}
\usepackage[T1]{fontenc}
\usepackage[spanish]{babel}
\usepackage[margin=2.3cm]{geometry}
\usepackage{amsmath}
\usepackage{booktabs}
\usepackage{float}
\usepackage{array}
\usepackage{enumitem}
\usepackage{tikz}
\usetikzlibrary{shapes.geometric, arrows.meta, positioning}
\usepackage[hidelinks]{hyperref}

\renewcommand{\contentsname}{Índice}
\renewcommand{\figurename}{Figura}
\renewcommand{\tablename}{Tabla}

\setlist{itemsep=2pt, topsep=4pt}

\tikzset{
  bloque/.style={rectangle, rounded corners, draw=black, thick, fill=blue!8,
                 text width=7.5cm, align=center, font=\small, inner sep=6pt},
  bloqueMCU/.style={bloque, fill=green!12},
  bloqueSPI/.style={bloque, fill=orange!25},
  bloqueSBC/.style={bloque, fill=violet!12},
  flujo/.style={-{Stealth[length=2.5mm]}, thick},
  bidir/.style={{Stealth[length=2.5mm]}-{Stealth[length=2.5mm]}, thick},
  inicio/.style={rectangle, rounded corners=6pt, draw=black, thick, fill=green!15,
                 minimum width=3.5cm, minimum height=0.8cm, align=center, font=\small},
  proc/.style={rectangle, draw=black, thick, fill=blue!8, text width=5.2cm,
               align=center, font=\small, minimum height=0.9cm, inner sep=4pt},
  dec/.style={diamond, aspect=2.4, draw=black, thick, fill=yellow!20,
              text width=3.0cm, align=center, font=\footnotesize, inner sep=1pt},
}

% Datos del proyecto (un solo lugar para mantenerlos)
\newcommand{\Universidad}{Universidad Nacional de Colombia --- Sede Manizales}
\newcommand{\Facultad}{Facultad de Ingeniería y Arquitectura}
\newcommand{\Carrera}{Carrera de Ingeniería Electrónica}
\newcommand{\Materia}{Programación de sistemas Linux embebido}
\newcommand{\TituloProyecto}{Sistema Híbrido de Adquisición Analógica mediante SPI\\[0.3cm]y Dashboard Embebido en Linux}
\newcommand{\FechaDoc}{Manizales, noviembre de 2026}

% Comando de color-neutral para valores a completar por el equipo
\newcommand{\completar}[1]{\textit{[#1]}}
% ============================================================
% Portada común: \portada{Título del documento}
% ============================================================
\newcommand{\portada}[1]{%
\begin{titlepage}
\centering
{\large \Universidad\par}
{\large \Facultad\par}
{\large \Carrera\par}
\vspace{1.2cm}
{\large Materia: \textbf{\Materia}\par}
\vspace{1.8cm}
{\LARGE\bfseries \TituloProyecto\par}
\vspace{0.5cm}
{\Large #1\par}
\vspace{1.5cm}

\small
\begin{tabular}{@{}p{0.5cm} p{4.2cm} p{2.8cm} p{6.0cm}@{}}
\toprule
& \textbf{Integrante} & \textbf{C.C.} & \textbf{Rol técnico} \\
\midrule
1. & \hrulefill & \hrulefill & Ingeniero de Documentación Técnica y Gestión de Configuración \\[0.7cm]
2. & \hrulefill & \hrulefill & Arquitecto de Sistemas Embebidos \\[0.7cm]
3. & \hrulefill & \hrulefill & Ingeniero de Software y Firmware Embebido \\[0.7cm]
4. & \hrulefill & \hrulefill & Ingeniero de Verificación y Validación (V\&V) \\[0.7cm]
\bottomrule
\end{tabular}

\vfill
\FechaDoc
\end{titlepage}
}

\begin{document}

\portada{Documento 7: Guía de operación y pruebas manuales de funcionamiento}

\tableofcontents
\newpage

% ============================================================
\section{Objetivo}

Esta guía permite que \textbf{una persona}, sin ayuda de la IA, ponga en marcha el osciloscopio desde las conexiones hasta la página web y compruebe su funcionamiento con pruebas sencillas. Las pruebas no requieren instrumentos de laboratorio: la propia Raspberry Pi genera una onda cuadrada de frecuencia conocida. Si hay un generador de funciones disponible, se indica cómo usarlo.

Cada prueba tiene un resultado esperado y un criterio de aceptación. Los resultados se anotan en la hoja de registro de la Sección~\ref{sec:hoja}.

\section{Materiales}

\begin{itemize}
  \item ESP32-S3-DevKitC-1 con el firmware cargado, Raspberry Pi 4 con el software instalado y los cables SPI conectados (Tabla~\ref{tab:spi}).
  \item Cable micro-USB para alimentar el ESP32-S3 (puerto «UART» de la placa) y fuente de la Raspberry Pi.
  \item PC en la misma red que la Raspberry Pi, con navegador web.
  \item Cables tipo \emph{jumper} macho-hembra.
  \item Dos resistencias iguales de 10 k$\Omega$ (para un divisor de tensión) y, si se dispone, un potenciómetro de 10 k$\Omega$ y una resistencia de 1 k$\Omega$.
  \item Cronómetro (el del teléfono sirve).
  \item Opcional: generador de funciones.
\end{itemize}

\paragraph{Reglas de seguridad.}
\begin{itemize}
  \item Las entradas CH1 (GPIO1) y CH2 (GPIO2) admiten \textbf{solo de 0 a 3{,}3 V}. Nunca conectar 5 V.
  \item No unir entre placas los pines 5V ni 3V3; cada placa tiene su propia alimentación. Solo se unen las tierras (GND).
  \item Hacer los cambios de cableado de las entradas con las placas encendidas está permitido, pero con cuidado de no tocar pines vecinos.
\end{itemize}

% ============================================================
\section{Puesta en marcha paso a paso}

\subsection{Paso 1: conexiones}

\begin{figure}[H]
\centering
\includegraphics[width=\textwidth]{Circuito.png}
\caption{Circuito de referencia (Fritzing): SPI entre la Raspberry Pi y el ESP32-S3, y adaptador CP2102 para la consola UART (paso 4b). La placa real es una ESP32-S3-DevKitC-1; las entradas CH1 y CH2 se conectan según cada prueba.}
\label{fig:circuito7}
\end{figure}

Con la DevKitC-1 vista con los conectores USB hacia abajo, los pines se cuentan desde arriba.

\begin{table}[H]
\centering
\small
\begin{tabular}{@{}p{1.6cm} p{2.0cm} p{4.4cm} p{4.4cm}@{}}
\toprule
\textbf{Señal} & \textbf{ESP32-S3} & \textbf{Posición en la DevKitC-1} & \textbf{Raspberry Pi 4 (pin físico)} \\
\midrule
CS   & GPIO10 & Cabecera izquierda, pin 16 & Pin 24 (GPIO8 / CE0) \\
MOSI & GPIO11 & Cabecera izquierda, pin 17 & Pin 19 (GPIO10 / MOSI) \\
SCLK & GPIO12 & Cabecera izquierda, pin 18 & Pin 23 (GPIO11 / SCLK) \\
MISO & GPIO13 & Cabecera izquierda, pin 19 & Pin 21 (GPIO9 / MISO) \\
GND  & GND    & Cabecera izquierda, pin 22 & Pin 25 (GND) \\
\bottomrule
\end{tabular}
\caption{Enlace SPI.}
\label{tab:spi}
\end{table}

Las entradas analógicas son \textbf{CH1 = GPIO1} (cabecera derecha, pin 4) y \textbf{CH2 = GPIO2} (cabecera derecha, pin 5). Cada prueba indica qué conectar en ellas.

\subsection{Paso 2: encendido}

\begin{enumerate}
  \item Encender la Raspberry Pi y esperar alrededor de 1 minuto a que arranque.
  \item Conectar el ESP32-S3 por el micro-USB marcado «UART» (al PC o a cualquier cargador USB).
\end{enumerate}

No hace falta ningún comando: el servicio \texttt{oscd} y el servidor Nginx arrancan solos con la Raspberry Pi, y el firmware arranca solo con el ESP32-S3. El orden de encendido no importa.

\subsection{Paso 3: abrir el dashboard}

En el navegador del PC abrir:

\begin{verbatim}
http://10.138.40.28/
\end{verbatim}

Si la dirección de la Raspberry Pi cambió, averiguarla por la consola UART (paso 4b) o en el router; en la Raspberry Pi se obtiene con \texttt{hostname -I}.

El indicador superior debe decir \textbf{«en vivo»} y la tabla inferior debe mostrar el contador de ventanas recibidas aumentando y \textbf{0 errores de CRC}.

\subsection{Paso 4 (opcional): verificar el servicio en la Raspberry Pi}

Desde una terminal del PC (PowerShell o WSL):

\begin{verbatim}
ssh juaro@10.138.40.28
systemctl status oscd nginx --no-pager
curl -s http://localhost/data/window.json | head -c 300; echo
journalctl -u oscd -n 20 --no-pager
\end{verbatim}

Los dos servicios deben aparecer \texttt{active (running)}. Si \texttt{oscd} no está activo:

\begin{verbatim}
sudo systemctl restart oscd
\end{verbatim}

\subsection{Paso 4b (alternativa a SSH): consola de la Raspberry Pi por UART}

Sirve cuando no hay red o no se conoce la dirección IP de la Raspberry Pi. Se usa un adaptador USB-UART CP2102 conectado así:

\begin{table}[H]
\centering
\small
\begin{tabular}{@{}p{3.6cm} p{6.0cm}@{}}
\toprule
\textbf{Adaptador CP2102} & \textbf{Raspberry Pi 4 (pin físico)} \\
\midrule
TX  & Pin 10 (GPIO15 / RXD) \\
RX  & Pin 8 (GPIO14 / TXD) \\
GND & Pin 6 (GND) \\
\bottomrule
\end{tabular}
\caption{Cableado de la consola UART. TX y RX se cruzan; no conectar los pines 5V ni 3V3 del adaptador.}
\end{table}

\paragraph{Opción A: desde WSL (no requiere instalar nada).}

\begin{enumerate}
  \item En PowerShell (no requiere administrador), identificar el adaptador y pasarlo a WSL. El BUSID es el de la línea \texttt{CP2102 USB to UART Bridge Controller}; si dice \emph{Not shared}, ejecutar una sola vez \texttt{usbipd bind --busid <BUSID>} en una PowerShell de administrador.
\begin{verbatim}
usbipd list
usbipd attach --wsl --busid <BUSID>
\end{verbatim}
  \item En Ubuntu (WSL), abrir la terminal serie:
\begin{verbatim}
ls /dev/ttyUSB*
picocom -b 115200 /dev/ttyUSB0
\end{verbatim}
  Si hay dos puertos, el de la Raspberry Pi es el que responde con \texttt{login:}; el otro es el ESP32-S3 (imprime líneas \texttt{STAT}).
  \item Presionar \textbf{Enter}. Aparece:
\begin{verbatim}
Juaromero login:
\end{verbatim}
  Escribir el usuario \texttt{juaro} y la contraseña (no se ve al escribirla). Desde aquí se pueden ejecutar los mismos comandos del paso 4, por ejemplo:
\begin{verbatim}
hostname -I
systemctl status oscd nginx --no-pager
\end{verbatim}
  \item Para terminar: \texttt{exit} (cierra la sesión de la Raspberry Pi) y luego \textbf{Ctrl+A} seguido de \textbf{Ctrl+X} (cierra picocom).
  \item Devolver el adaptador a Windows:
\begin{verbatim}
usbipd detach --busid <BUSID>
\end{verbatim}
\end{enumerate}

Con \texttt{screen} en lugar de picocom: \texttt{screen /dev/ttyUSB0 115200}; se sale con \textbf{Ctrl+A}, \textbf{K}, \textbf{Y}. No salir con \textbf{Ctrl+A}, \textbf{D}, porque eso deja la sesión abierta en segundo plano y bloquea el puerto.

\paragraph{Opción B: desde Windows con PuTTY.} Instalar el controlador \emph{CP210x Universal Windows Driver} de Silicon Labs; reconectar el adaptador y anotar en el Administrador de dispositivos el puerto \texttt{COMx} asignado. En PuTTY elegir \emph{Serial}, ese puerto y velocidad 115200, y presionar Enter para ver el \texttt{login:}. El adaptador no puede estar adjunto a WSL al mismo tiempo.

\paragraph{Configuración de la Raspberry Pi (ya aplicada).} La consola UART funciona porque en \texttt{/boot/firmware/config.txt} están \texttt{enable\_uart=1} y \texttt{dtoverlay=disable-bt}, y porque el servicio \texttt{serial-getty@ttyAMA0} está habilitado. Para desactivarla:
\begin{verbatim}
sudo systemctl disable --now serial-getty@ttyAMA0
\end{verbatim}
y para volver a activarla:
\begin{verbatim}
sudo systemctl enable --now serial-getty@ttyAMA0
\end{verbatim}

\subsection{Paso 5 (opcional): monitor serie del ESP32-S3}

El monitor serie muestra cada segundo una línea \texttt{STAT} y la última ventana. En PowerShell (no requiere administrador):

\begin{verbatim}
usbipd list
usbipd attach --wsl --busid <BUSID del CP2102N>
\end{verbatim}

En Ubuntu (WSL):

\begin{verbatim}
ls /dev/ttyUSB*
picocom -b 115200 /dev/ttyUSB1
\end{verbatim}

Si hay dos puertos, el del ESP32-S3 es el que responde con líneas \texttt{STAT}. Para salir de picocom: \textbf{Ctrl+A} y luego \textbf{Ctrl+X}. Las teclas 0, 1, 2 y 3 cambian el preset desde el monitor serie.

\subsection{Paso 6 (solo si se modificó el firmware): volver a programar el ESP32-S3}

Con la placa adjunta a WSL (paso 5) y el monitor serie cerrado:

\begin{verbatim}
source ~/zephyrproject/.venv/bin/activate
cd ~/zephyrproject
R=/mnt/c/Users/juaro/OneDrive/Desktop
P=$R/Programacion_de_Sistemas_Linux_Embebidos
APP=$P/Proyecto_2/Codigo/Firmware_ESP32S3
west build -p auto -b esp32s3_devkitc/esp32s3/procpu "$APP" \
  -d ~/build/osc
west flash -d ~/build/osc --esp-device /dev/ttyUSB1
\end{verbatim}

\subsection{Paso 7: apagado}

\begin{verbatim}
ssh juaro@10.138.40.28 sudo shutdown -h now
\end{verbatim}

Esperar a que el LED verde de la Raspberry Pi deje de parpadear antes de quitarle la alimentación. El ESP32-S3 se puede desconectar en cualquier momento.

% ============================================================
\section{Generar una señal de prueba con la Raspberry Pi}

La Raspberry Pi puede producir una onda cuadrada de 0 a 3{,}3 V en su pin físico \textbf{12} (GPIO18) con el comando \texttt{gpioset}. Conexión: pin 12 de la Raspberry Pi $\rightarrow$ (resistencia de 1 k$\Omega$, recomendada) $\rightarrow$ GPIO1 (CH1) del ESP32-S3. La tierra ya es común por el cable GND del SPI.

Desde una sesión SSH en la Raspberry Pi:

\begin{verbatim}
gpioset -c gpiochip0 -t 8ms GPIO18=1
\end{verbatim}

El pin cambia de estado cada 8 ms, de modo que el período es 16 ms ($f = 62{,}5$ Hz). Se detiene con \textbf{Ctrl+C}. Frecuencia de la onda: $f = 1 / (2\,t)$, donde $t$ es el valor de \texttt{-t}.

\paragraph{Limitación.} \texttt{gpioset} es un programa de Linux, no un generador de precisión: los flancos tienen variaciones del orden de decenas de microsegundos. Sirve para pruebas de funcionamiento hasta $\approx$ 1 kHz. Para medir con exactitud se usa un generador de funciones (Sección~\ref{sec:generador}).

\begin{table}[H]
\centering
\small
\begin{tabular}{@{}p{2.4cm} p{2.4cm} p{3.0cm} p{2.4cm} p{3.6cm}@{}}
\toprule
\textbf{Preset} & \textbf{Ventana} & \textbf{Comando (\texttt{-t})} & \textbf{Frecuencia} & \textbf{Períodos esperados en pantalla} \\
\midrule
100 Sa/s     & 640 ms     & \texttt{-t 80ms}  & 6{,}25 Hz & 4 \\
1\,000 Sa/s  & 64 ms      & \texttt{-t 8ms}   & 62{,}5 Hz & 4 \\
10\,000 Sa/s & 6{,}4 ms   & \texttt{-t 800us} & 625 Hz    & 4 \\
41\,667 Sa/s & 1{,}536 ms & \texttt{-t 500us} & 1 kHz     & $\approx$ 1{,}5 \\
\bottomrule
\end{tabular}
\caption{Señal de prueba recomendada para cada preset.}
\label{tab:senal}
\end{table}

% ============================================================
\section{Pruebas de funcionamiento}

\subsection*{P-01. Arranque y dashboard en vivo}
\begin{itemize}
  \item \textbf{Pasos:} seguir los pasos 1 a 3. Observar la página durante 1 minuto.
  \item \textbf{Esperado:} indicador «en vivo»; el número de ventana aumenta; 0 errores de CRC.
  \item \textbf{Criterio:} la página se actualiza sin errores durante 1 minuto.
\end{itemize}

\subsection*{P-02. Nivel de tierra (0 V)}
\begin{itemize}
  \item \textbf{Montaje:} CH1 (GPIO1) y CH2 (GPIO2) conectados a un pin GND del ESP32-S3.
  \item \textbf{Esperado:} ambas trazas planas cerca de 0 V.
  \item \textbf{Criterio:} ambas trazas por debajo de 0{,}15 V.
\end{itemize}

\subsection*{P-03. Nivel alto (3{,}3 V)}
\begin{itemize}
  \item \textbf{Montaje:} CH1 y CH2 conectados al pin 3V3 del ESP32-S3.
  \item \textbf{Esperado:} ambas trazas planas arriba, en $\approx$ 3{,}1 V (es el fondo de escala del ADC; 3{,}3 V lo satura).
  \item \textbf{Criterio:} ambas trazas por encima de 2{,}9 V.
\end{itemize}

\subsection*{P-04. Canales no intercambiados}
\begin{itemize}
  \item \textbf{Montaje:} CH1 a 3V3 y CH2 a GND. Después, al revés.
  \item \textbf{Esperado:} en la leyenda, CH1 es la traza naranja y CH2 la azul; la traza alta debe corresponder al canal conectado a 3V3.
  \item \textbf{Criterio:} los canales siguen al cable en los dos casos.
\end{itemize}

\subsection*{P-05. Valor intermedio con divisor de tensión}
\begin{itemize}
  \item \textbf{Montaje:} dos resistencias de 10 k$\Omega$ en serie entre 3V3 y GND; el punto medio a CH1. CH2 a GND.
  \item \textbf{Esperado:} CH1 plana en $\approx$ 1{,}65 V.
  \item \textbf{Criterio:} CH1 entre 1{,}40 V y 1{,}90 V. El ADC del ESP32-S3 no está calibrado, por eso la tolerancia es amplia.
\end{itemize}

\subsection*{P-06. Barrido con potenciómetro (opcional)}
\begin{itemize}
  \item \textbf{Montaje:} potenciómetro entre 3V3 y GND, el cursor a CH1.
  \item \textbf{Pasos:} girar lentamente de un extremo al otro, con el refresco en 1 Hz.
  \item \textbf{Esperado:} la traza sube y baja de forma continua entre $\approx$ 0 V y $\approx$ 3{,}1 V, sin saltos.
  \item \textbf{Criterio:} respuesta continua y monótona.
\end{itemize}

\subsection*{P-07. Frecuencia de muestreo en los cuatro presets}
\begin{itemize}
  \item \textbf{Montaje:} onda cuadrada de la Raspberry Pi en CH1 (Tabla~\ref{tab:senal}); CH2 a GND.
  \item \textbf{Pasos:} para cada preset, pulsar su botón en el dashboard, ejecutar el comando \texttt{gpioset} correspondiente y contar los períodos completos en la pantalla. Leer también la duración de un período en el eje de tiempo (ms).
  \item \textbf{Esperado:} el número de períodos de la Tabla~\ref{tab:senal}; período de $2t$.
  \item \textbf{Criterio:} período medido dentro de $\pm 10\,\%$ del esperado. La tolerancia cubre la imprecisión de \texttt{gpioset}, no la del osciloscopio.
\end{itemize}

\subsection*{P-08. Cambio de preset desde el dashboard}
\begin{itemize}
  \item \textbf{Pasos:} con la onda de 62{,}5 Hz en CH1, pulsar en orden 1 kSa/s, 10 kSa/s, 41{,}7 kSa/s y 100 Sa/s.
  \item \textbf{Esperado:} el botón activo cambia, el encabezado muestra la nueva frecuencia de muestreo y el eje de tiempo cambia de escala (64 ms, 6{,}4 ms, 1{,}5 ms y 640 ms).
  \item \textbf{Criterio:} cada cambio se ve en 2 refrescos o menos.
\end{itemize}

\subsection*{P-09. Tasa de refresco}
\begin{itemize}
  \item \textbf{Pasos:} elegir «0,25 Hz» en el selector de refresco. Con la onda cuadrada aplicada, cronometrar 10 actualizaciones de la gráfica. Repetir con «1 Hz».
  \item \textbf{Esperado:} 40 s para 10 actualizaciones a 0{,}25 Hz y 10 s a 1 Hz.
  \item \textbf{Criterio:} dentro de $\pm 2$ s (0{,}25 Hz) y $\pm 1$ s (1 Hz). La tolerancia es la de un cronómetro manual.
\end{itemize}

\subsection*{P-10. Submuestreo (aliasing) --- demostración}
\begin{itemize}
  \item \textbf{Montaje:} onda de 62{,}5 Hz (\texttt{-t 8ms}) en CH1, con el preset de 100 Sa/s.
  \item \textbf{Esperado:} la señal aparece con una frecuencia \emph{falsa} de $\approx$ 37{,}5 Hz, porque 62{,}5 Hz supera la frecuencia de Nyquist (50 Hz). Con el preset de 1 kSa/s se ve la frecuencia correcta.
  \item \textbf{Criterio:} se observa la frecuencia aparente distinta a 100 Sa/s y la correcta a 1 kSa/s. Demuestra el teorema de muestreo.
\end{itemize}

\subsection*{P-11. Coincidencia con el monitor serie}
\begin{itemize}
  \item \textbf{Pasos:} con CH1 en el divisor (P-05), abrir el monitor serie (paso 5) y comparar los valores de las líneas \texttt{WIN} con la gráfica. El dashboard muestra voltios; el monitor serie muestra códigos de 0 a 4095 (voltios $\approx$ código $\times 3{,}1 / 4095$).
  \item \textbf{Criterio:} los dos muestran el mismo nivel.
\end{itemize}

\subsection*{P-12. Recuperación ante fallos}
\begin{itemize}
  \item \textbf{a)} Desconectar el USB del ESP32-S3 durante 10 s y volver a conectarlo. \textbf{Esperado:} el indicador pasa a «sin ventana nueva» y vuelve a «en vivo» por sí solo, con el mismo preset que estaba elegido.
  \item \textbf{b)} En la Raspberry Pi: \texttt{sudo systemctl restart oscd}. \textbf{Esperado:} el dashboard se recupera en menos de 5 s.
  \item \textbf{c)} Reiniciar la Raspberry Pi con \texttt{sudo reboot}. \textbf{Esperado:} tras $\approx$ 1 minuto la página vuelve a funcionar sin ejecutar ningún comando.
  \item \textbf{Criterio:} recuperación sin intervención en los tres casos.
\end{itemize}

\subsection*{P-13. Estabilidad}
\begin{itemize}
  \item \textbf{Pasos:} con el preset de 41{,}7 kSa/s y la onda de 1 kHz, dejar el dashboard abierto 10 minutos.
  \item \textbf{Criterio:} la página sigue «en vivo» y los errores de CRC permanecen en 0.
\end{itemize}

\subsection*{P-14. Mediciones automáticas del dashboard}
\begin{itemize}
  \item \textbf{Montaje:} onda cuadrada de 62{,}5 Hz (\texttt{-t 8ms}) en CH1, preset de 1 kSa/s.
  \item \textbf{Esperado:} en la tarjeta de CH1, frecuencia $\approx$ 62{,}5 Hz, mínimo $\approx$ 0 V, máximo $\approx$ 3{,}1 V y el aviso verde «16 muestras por período». Con el preset de 100 Sa/s el aviso cambia a rojo (pocas muestras por período) y recomienda subir la tasa.
  \item \textbf{Criterio:} frecuencia dentro de $\pm 10\,\%$ y recomendación coherente con la Tabla~\ref{tab:senal}.
\end{itemize}

\subsection*{P-15. Tensiones negativas (requiere el acondicionador)}
\begin{itemize}
  \item \textbf{Montaje:} acondicionador de 18 k$\Omega$ / 10 k$\Omega$ / 18 k$\Omega$ (Documento 1) entre la señal y CH1. En el dashboard, CH1 en modo «$\pm 5$ V con acondicionador».
  \item \textbf{Pasos:} (1) unir la entrada del acondicionador a GND y pulsar «Calibrar cero» en CH1; (2) conectar la entrada a $+3{,}3$ V (pin 3V3); (3) si hay fuente o generador bipolar, aplicar $-3$ V o una senoide de $\pm 3$ V.
  \item \textbf{Esperado:} (1) CH1 en $\approx$ 0 V; (2) $\approx +3{,}3$ V; (3) $\approx -3$ V o la senoide centrada en 0 V con mínimo $\approx -3$ V.
  \item \textbf{Criterio:} lecturas dentro de $\pm 0{,}3$ V. \textbf{No} conectar nunca una tensión negativa directamente a GPIO1 o GPIO2, sin el acondicionador.
\end{itemize}

\subsection*{P-16. Exportar CSV y pausa}
\begin{itemize}
  \item \textbf{Pasos:} pulsar «Pausar» (la gráfica se congela y el indicador dice «en pausa»), pulsar «Exportar CSV» y abrir el archivo en una hoja de cálculo; pulsar «Reanudar».
  \item \textbf{Criterio:} el archivo tiene 64 filas con tiempo, tensiones y códigos de ambos canales, y la gráfica se reanuda.
\end{itemize}

% ============================================================
\section{Con generador de funciones (opcional)}
\label{sec:generador}

Configurar el generador en senoidal, con \textbf{offset de 1{,}5 V y amplitud de 1 V pico a pico} (la señal queda entre 1{,}0 V y 2{,}0 V, dentro del rango del ADC). Comprobar con el propio generador o con un multímetro que nunca supere 3{,}3 V \emph{antes} de conectarlo. Su tierra se une a un GND del ESP32-S3 y la señal a CH1.

Con las frecuencias de la Tabla~\ref{tab:senal} (o, para el preset de 41{,}7 kSa/s, 2{,}6 kHz, que da 4 períodos) se repite P-07 midiendo el período en la pantalla. Con un generador el criterio puede ser más estricto: $\pm 2\,\%$. Aplicando la misma señal a CH1 y CH2 se ve además el pequeño desfase entre canales (CH2 se muestrea un intervalo de conversión después de CH1).

% ============================================================
\section{Hoja de registro}
\label{sec:hoja}

\noindent Responsable de las pruebas: \hrulefill \quad Fecha: \rule{3cm}{0.4pt}

\begin{table}[H]
\centering
\small
\begin{tabular}{@{}p{1.2cm} p{4.4cm} p{5.6cm} p{1.4cm}@{}}
\toprule
\textbf{Prueba} & \textbf{Criterio} & \textbf{Resultado observado} & \textbf{¿Cumple?} \\
\midrule
P-01 & En vivo 1 min, CRC = 0 & & \\[0.5cm]
P-02 & CH1 y CH2 $<$ 0{,}15 V & CH1: \qquad CH2: & \\[0.5cm]
P-03 & CH1 y CH2 $>$ 2{,}9 V & CH1: \qquad CH2: & \\[0.5cm]
P-04 & Canales siguen al cable & & \\[0.5cm]
P-05 & CH1 entre 1{,}40 y 1{,}90 V & CH1: & \\[0.5cm]
P-06 & Respuesta continua & & \\[0.5cm]
P-07 & Período $\pm 10\,\%$ & 100: \quad 1k: \quad 10k: \quad 41k: & \\[0.5cm]
P-08 & Cambio en $\leq$ 2 refrescos & & \\[0.5cm]
P-09 & 40 s $\pm$ 2 s / 10 s $\pm$ 1 s & 0{,}25 Hz: \quad 1 Hz: & \\[0.5cm]
P-10 & Alias visible a 100 Sa/s & Frecuencia aparente: & \\[0.5cm]
P-11 & Serie = dashboard & & \\[0.5cm]
P-12 & Recupera sin intervención & a) \quad b) \quad c) & \\[0.5cm]
P-13 & 10 min, CRC = 0 & CRC final: & \\[0.5cm]
P-14 & Frecuencia $\pm 10\,\%$, aviso coherente & Frecuencia: & \\[0.5cm]
P-15 & $\pm 0{,}3$ V (con acondicionador) & 0 V: \quad +3,3 V: \quad $-3$ V: & \\[0.5cm]
P-16 & CSV de 64 filas; pausa y reanudación & & \\[0.5cm]
\bottomrule
\end{tabular}
\caption{Registro de las pruebas manuales.}
\end{table}

\noindent Observaciones: \hrulefill

\vspace{0.6cm}
\noindent\hrulefill

\vspace{1.2cm}
\noindent Firma: \rule{6cm}{0.4pt}

% ============================================================
\section{Solución de problemas}

\begin{table}[H]
\centering
\small
\begin{tabular}{@{}p{4.6cm} p{4.6cm} p{5.4cm}@{}}
\toprule
\textbf{Síntoma} & \textbf{Causa probable} & \textbf{Qué hacer} \\
\midrule
La página no abre & Raspberry Pi apagada, otra red u otra dirección IP & Verificar la red; obtener la IP con \texttt{hostname -I} por UART \\
«sin datos» en la página & \texttt{oscd} detenido & \texttt{sudo systemctl restart oscd} \\
«sin ventana nueva» permanente & ESP32-S3 sin alimentación o cable SPI suelto & Revisar USB y los cinco cables SPI \\
Errores de CRC aumentando & Cables SPI largos o flojos & Acortar y asegurar cables; revisar GND \\
Trazas con ruido de 0 a 3 V & Entrada sin conectar (flotante) & Es normal: conectar la señal o GND \\
La traza alta no corresponde al canal conectado & Cable en el pin equivocado & Revisar GPIO1 / GPIO2 en la cabecera derecha \\
\texttt{Device or resource busy} en picocom & Otra terminal abierta en el puerto & \texttt{screen -r} o cerrar la otra terminal \\
\texttt{gpioset}: \emph{line busy} & Otro \texttt{gpioset} sigue activo & Detenerlo con Ctrl+C o \texttt{pkill gpioset} \\
Consola UART sin \texttt{login:} & TX y RX sin cruzar, falta GND o velocidad distinta de 115200 & Revisar el cableado del paso 4b y la velocidad \\
CP2102 con error en Windows & Falta el controlador CP210x & Instalarlo o usar la opción A (WSL) \\
\bottomrule
\end{tabular}
\caption{Problemas frecuentes.}
\end{table}

\end{document}
